%%%%%%%%%%%%%%%%%%%%%%%%%%%%%%%%%%%%%%%%%%%%%%%%%%%%%%%%%%%%
% 10.9 MULTIOBJECTIVE OPTIMIZATION
% The Composition of Advanced Mathematics
%%%%%%%%%%%%%%%%%%%%%%%%%%%%%%%%%%%%%%%%%%%%%%%%%%%%%%%%%%%%

\section{Multiobjective Optimization}
\label{sec:multiobjective-optimization}

\prerequisite{
Optimization fundamentals, convex optimization, nonlinear programming,
linear programming, vectors, inequalities, feasible sets, and basic
decision theory.
}

\learningobjectives{
By the end of this section, the reader should be able to:
\begin{itemize}
    \item define a multiobjective optimization problem;
    \item distinguish single-objective and multiobjective optimization;
    \item identify conflicting objectives;
    \item define dominance and Pareto dominance;
    \item define Pareto-optimal and weakly Pareto-optimal solutions;
    \item distinguish decision space and objective space;
    \item define the Pareto set and Pareto frontier;
    \item recognize dominated and nondominated solutions;
    \item understand tradeoffs among objectives;
    \item define ideal and reference objective vectors;
    \item understand the nadir point conceptually;
    \item normalize objectives with different scales;
    \item formulate weighted-sum scalarizations;
    \item understand limitations of weighted sums;
    \item understand supported and unsupported Pareto points;
    \item formulate epsilon-constraint problems;
    \item formulate goal-programming models;
    \item understand lexicographic optimization;
    \item formulate reference-point and achievement scalarizing models;
    \item recognize compromise programming;
    \item understand Chebyshev scalarization;
    \item formulate multiobjective linear programs;
    \item formulate multiobjective convex programs;
    \item understand first-order Pareto optimality conceptually;
    \item understand weighted KKT conditions conceptually;
    \item recognize proper efficiency conceptually;
    \item understand decision-maker preferences;
    \item distinguish a posteriori, a priori, and interactive methods;
    \item understand generation of Pareto frontiers;
    \item recognize evolutionary multiobjective optimization conceptually;
    \item understand hypervolume and diversity conceptually;
    \item identify sensitivity of solutions to objective weights;
    \item understand fairness and competing stakeholder objectives;
    \item connect multiobjective optimization with engineering,
          economics, finance, planning, sustainability, and operations
          research.
\end{itemize}
}

%%%%%%%%%%%%%%%%%%%%%%%%%%%%%%%%%%%%%%%%%%%%%%%%%%%%%%%%%%%%
\subsection{What Is Multiobjective Optimization?}
\label{subsec:what-is-multiobjective-optimization}
%%%%%%%%%%%%%%%%%%%%%%%%%%%%%%%%%%%%%%%%%%%%%%%%%%%%%%%%%%%%

A multiobjective optimization problem has more than one objective function
to optimize simultaneously.

A general form is

\[
\boxed{
\min_{\mathbf{x}\in\mathcal{F}}
\begin{pmatrix}
f_1(\mathbf{x})\\
f_2(\mathbf{x})\\
\vdots\\
f_k(\mathbf{x})
\end{pmatrix}.
}
\]

Here,

\[
\mathbf{x}
\]

is the decision vector,

\[
\mathcal{F}
\]

is the feasible set, and

\[
f_1,\ldots,f_k
\]

are the objectives.

%%%%%%%%%%%%%%%%%%%%%%%%%%%%%%%%%%%%%%%%%%%%%%%%%%%%%%%%%%%%
\subsection{Conflicting Objectives}
\label{subsec:conflicting-objectives}
%%%%%%%%%%%%%%%%%%%%%%%%%%%%%%%%%%%%%%%%%%%%%%%%%%%%%%%%%%%%

If every objective were minimized by the same feasible point, the problem
would effectively reduce to a single decision.

The interesting case occurs when objectives conflict.

For example, one may want to:

\[
\boxed{
\text{minimize cost}
}
\]

while also

\[
\boxed{
\text{maximize quality}.
}
\]

Improving one objective may worsen another.

%%%%%%%%%%%%%%%%%%%%%%%%%%%%%%%%%%%%%%%%%%%%%%%%%%%%%%%%%%%%
\subsection{Examples of Multiple Objectives}
\label{subsec:examples-multiple-objectives}
%%%%%%%%%%%%%%%%%%%%%%%%%%%%%%%%%%%%%%%%%%%%%%%%%%%%%%%%%%%%

Common competing objectives include:

\[
\boxed{
\text{cost}
}
\quad\text{versus}\quad
\boxed{
\text{performance},
}
\]

\[
\boxed{
\text{return}
}
\quad\text{versus}\quad
\boxed{
\text{risk},
}
\]

\[
\boxed{
\text{speed}
}
\quad\text{versus}\quad
\boxed{
\text{energy use},
}
\]

and

\[
\boxed{
\text{profit}
}
\quad\text{versus}\quad
\boxed{
\text{environmental impact}.
}
\]

%%%%%%%%%%%%%%%%%%%%%%%%%%%%%%%%%%%%%%%%%%%%%%%%%%%%%%%%%%%%
\subsection{Vector-Valued Objective}
\label{subsec:vector-valued-objective}
%%%%%%%%%%%%%%%%%%%%%%%%%%%%%%%%%%%%%%%%%%%%%%%%%%%%%%%%%%%%

Define the objective vector

\[
\boxed{
\mathbf{f}(\mathbf{x})
=
\begin{pmatrix}
f_1(\mathbf{x})\\
\vdots\\
f_k(\mathbf{x})
\end{pmatrix}.
}
\]

Unlike ordinary scalar optimization, vectors do not generally possess a
natural total ordering.

For example,

\[
(2,7)
\]

and

\[
(4,3)
\]

cannot be ranked without additional preference information if both
components are to be minimized.

%%%%%%%%%%%%%%%%%%%%%%%%%%%%%%%%%%%%%%%%%%%%%%%%%%%%%%%%%%%%
\subsection{Decision Space and Objective Space}
\label{subsec:decision-objective-space}
%%%%%%%%%%%%%%%%%%%%%%%%%%%%%%%%%%%%%%%%%%%%%%%%%%%%%%%%%%%%

The decision vector

\[
\mathbf{x}
\]

belongs to decision space.

Its image

\[
\mathbf{f}(\mathbf{x})
\]

belongs to objective space.

Thus,

\[
\boxed{
\mathbf{x}
\in
\mathcal{F}
}
\]

is a feasible decision, while

\[
\boxed{
\mathbf{f}(\mathbf{x})
}
\]

is its objective vector.

The set

\[
\boxed{
\mathcal{Y}
=
\{
\mathbf{f}(\mathbf{x}):
\mathbf{x}\in\mathcal{F}
\}
}
\]

is the attainable objective set.

%%%%%%%%%%%%%%%%%%%%%%%%%%%%%%%%%%%%%%%%%%%%%%%%%%%%%%%%%%%%
\subsection{Pareto Dominance}
\label{subsec:pareto-dominance}
%%%%%%%%%%%%%%%%%%%%%%%%%%%%%%%%%%%%%%%%%%%%%%%%%%%%%%%%%%%%

Assume all objectives are minimized.

A feasible point \(\mathbf{x}\) dominates another feasible point
\(\mathbf{y}\) if

\[
\boxed{
f_i(\mathbf{x})
\leq
f_i(\mathbf{y})
\qquad
\text{for every }i,
}
\]

and

\[
\boxed{
f_j(\mathbf{x})
<
f_j(\mathbf{y})
}
\]

for at least one objective \(j\).

Thus \(\mathbf{x}\) is no worse in every objective and strictly better in
at least one.

%%%%%%%%%%%%%%%%%%%%%%%%%%%%%%%%%%%%%%%%%%%%%%%%%%%%%%%%%%%%
\subsection{Dominated Solutions}
\label{subsec:dominated-solutions}
%%%%%%%%%%%%%%%%%%%%%%%%%%%%%%%%%%%%%%%%%%%%%%%%%%%%%%%%%%%%

A feasible solution is dominated if another feasible solution improves at
least one objective without worsening any other objective.

Dominated solutions are usually unattractive because one can move to a
strictly better alternative.

%%%%%%%%%%%%%%%%%%%%%%%%%%%%%%%%%%%%%%%%%%%%%%%%%%%%%%%%%%%%
\subsection{Pareto-Optimal Solution}
\label{subsec:pareto-optimal-solution}
%%%%%%%%%%%%%%%%%%%%%%%%%%%%%%%%%%%%%%%%%%%%%%%%%%%%%%%%%%%%

\begin{definitionbox}{Pareto-Optimal Solution}
A feasible point \(\mathbf{x}^*\) is Pareto optimal if no feasible point
dominates it.
\end{definitionbox}

Equivalently, there is no feasible \(\mathbf{x}\) such that

\[
f_i(\mathbf{x})
\leq
f_i(\mathbf{x}^*)
\]

for every \(i\), with strict inequality for at least one objective.

%%%%%%%%%%%%%%%%%%%%%%%%%%%%%%%%%%%%%%%%%%%%%%%%%%%%%%%%%%%%
\subsection{Pareto Set}
\label{subsec:pareto-set}
%%%%%%%%%%%%%%%%%%%%%%%%%%%%%%%%%%%%%%%%%%%%%%%%%%%%%%%%%%%%

The Pareto set is the set of all Pareto-optimal decisions:

\[
\boxed{
\mathcal{P}
=
\{
\mathbf{x}\in\mathcal{F}:
\mathbf{x}
\text{ is Pareto optimal}
\}.
}
\]

This set lies in decision space.

%%%%%%%%%%%%%%%%%%%%%%%%%%%%%%%%%%%%%%%%%%%%%%%%%%%%%%%%%%%%
\subsection{Pareto Frontier}
\label{subsec:pareto-frontier}
%%%%%%%%%%%%%%%%%%%%%%%%%%%%%%%%%%%%%%%%%%%%%%%%%%%%%%%%%%%%

The Pareto frontier, also called the Pareto front, is the image of the
Pareto set in objective space:

\[
\boxed{
\mathcal{F}_P
=
\{
\mathbf{f}(\mathbf{x}):
\mathbf{x}\in\mathcal{P}
\}.
}
\]

Each point on the frontier represents a different tradeoff among
objectives.

%%%%%%%%%%%%%%%%%%%%%%%%%%%%%%%%%%%%%%%%%%%%%%%%%%%%%%%%%%%%
\subsection{Worked Example: Dominance}
\label{subsec:worked-dominance}
%%%%%%%%%%%%%%%%%%%%%%%%%%%%%%%%%%%%%%%%%%%%%%%%%%%%%%%%%%%%

\begin{workedexamplebox}{Comparing Objective Vectors}

Suppose both objectives are minimized.

Consider

\[
A=(4,7),
\]

\[
B=(5,9),
\]

and

\[
C=(3,10).
\]

Compare \(A\) and \(B\):

\[
4<5,
\]

and

\[
7<9.
\]

Therefore \(A\) dominates \(B\).

Compare \(A\) and \(C\):

\[
3<4
\]

for the first objective, but

\[
10>7
\]

for the second.

Neither dominates the other.

Thus,

\begin{answerbox}
\[
\boxed{
A\text{ dominates }B,
}
\]

while

\[
\boxed{
A\text{ and }C\text{ are incomparable by Pareto dominance}.
}
\]
\end{answerbox}

\end{workedexamplebox}

%%%%%%%%%%%%%%%%%%%%%%%%%%%%%%%%%%%%%%%%%%%%%%%%%%%%%%%%%%%%
\subsection{Weak Pareto Optimality}
\label{subsec:weak-pareto-optimality}
%%%%%%%%%%%%%%%%%%%%%%%%%%%%%%%%%%%%%%%%%%%%%%%%%%%%%%%%%%%%

A feasible point \(\mathbf{x}^*\) is weakly Pareto optimal if there is no
feasible point \(\mathbf{x}\) satisfying

\[
\boxed{
f_i(\mathbf{x})
<
f_i(\mathbf{x}^*)
}
\]

for every objective \(i\).

Thus no point can strictly improve all objectives simultaneously.

Every Pareto-optimal solution is weakly Pareto optimal, but the converse
need not hold.

%%%%%%%%%%%%%%%%%%%%%%%%%%%%%%%%%%%%%%%%%%%%%%%%%%%%%%%%%%%%
\subsection{Tradeoffs}
\label{subsec:multiobjective-tradeoffs}
%%%%%%%%%%%%%%%%%%%%%%%%%%%%%%%%%%%%%%%%%%%%%%%%%%%%%%%%%%%%

Moving along a Pareto frontier typically improves one objective while
worsening another.

For two objectives,

\[
f_1
\]

and

\[
f_2,
\]

a tradeoff may be interpreted as

\[
\boxed{
\text{amount of }f_2
\text{ sacrificed to improve }f_1.
}
\]

Multiobjective optimization does not remove such tradeoffs.

It makes them mathematically explicit.

%%%%%%%%%%%%%%%%%%%%%%%%%%%%%%%%%%%%%%%%%%%%%%%%%%%%%%%%%%%%
\subsection{Ideal Objective Vector}
\label{subsec:ideal-objective-vector}
%%%%%%%%%%%%%%%%%%%%%%%%%%%%%%%%%%%%%%%%%%%%%%%%%%%%%%%%%%%%

For minimization, define the ideal component

\[
\boxed{
f_i^{\mathrm{ideal}}
=
\min_{\mathbf{x}\in\mathcal{F}}
f_i(\mathbf{x}).
}
\]

The ideal objective vector is

\[
\boxed{
\mathbf{z}^{\mathrm{ideal}}
=
\begin{pmatrix}
f_1^{\mathrm{ideal}}\\
\vdots\\
f_k^{\mathrm{ideal}}
\end{pmatrix}.
}
\]

Each component is individually best possible.

%%%%%%%%%%%%%%%%%%%%%%%%%%%%%%%%%%%%%%%%%%%%%%%%%%%%%%%%%%%%
\subsection{Ideal Point May Be Unattainable}
\label{subsec:ideal-point-unattainable}
%%%%%%%%%%%%%%%%%%%%%%%%%%%%%%%%%%%%%%%%%%%%%%%%%%%%%%%%%%%%

The ideal point is constructed by minimizing each objective separately.

Those separate minima may occur at different decisions.

Therefore,

\[
\boxed{
\mathbf{z}^{\mathrm{ideal}}
}
\]

need not belong to the attainable objective set.

It is often used as a reference rather than an attainable solution.

%%%%%%%%%%%%%%%%%%%%%%%%%%%%%%%%%%%%%%%%%%%%%%%%%%%%%%%%%%%%
\subsection{Nadir Point}
\label{subsec:nadir-point}
%%%%%%%%%%%%%%%%%%%%%%%%%%%%%%%%%%%%%%%%%%%%%%%%%%%%%%%%%%%%

The nadir objective vector is conceptually formed from the worst objective
values occurring among Pareto-optimal solutions.

For minimization,

\[
\boxed{
z_i^{\mathrm{nadir}}
=
\max_{\mathbf{x}\in\mathcal{P}}
f_i(\mathbf{x}).
}
\]

The nadir point is often difficult to calculate exactly.

It is useful for scaling and visualization.

%%%%%%%%%%%%%%%%%%%%%%%%%%%%%%%%%%%%%%%%%%%%%%%%%%%%%%%%%%%%
\subsection{Utopia Point}
\label{subsec:utopia-point}
%%%%%%%%%%%%%%%%%%%%%%%%%%%%%%%%%%%%%%%%%%%%%%%%%%%%%%%%%%%%

Sometimes a reference point slightly better than the ideal point is used.

For minimization,

\[
\boxed{
z_i^{\mathrm{utopia}}
<
z_i^{\mathrm{ideal}}.
}
\]

This point is intentionally unattainable and can be useful in certain
scalarization methods.

%%%%%%%%%%%%%%%%%%%%%%%%%%%%%%%%%%%%%%%%%%%%%%%%%%%%%%%%%%%%
\subsection{Normalization}
\label{subsec:multiobjective-normalization}
%%%%%%%%%%%%%%%%%%%%%%%%%%%%%%%%%%%%%%%%%%%%%%%%%%%%%%%%%%%%

Suppose one objective has values around

\[
10^6,
\]

while another has values around

\[
10.
\]

Directly weighting these objectives can produce misleading results.

A normalized objective may be

\[
\boxed{
\widetilde{f}_i(\mathbf{x})
=
\frac{
f_i(\mathbf{x})
-
z_i^{\mathrm{ideal}}
}{
z_i^{\mathrm{nadir}}
-
z_i^{\mathrm{ideal}}
}.
}
\]

This places objectives on comparable scales when suitable estimates are
available.

%%%%%%%%%%%%%%%%%%%%%%%%%%%%%%%%%%%%%%%%%%%%%%%%%%%%%%%%%%%%
\subsection{Scalarization}
\label{subsec:multiobjective-scalarization}
%%%%%%%%%%%%%%%%%%%%%%%%%%%%%%%%%%%%%%%%%%%%%%%%%%%%%%%%%%%%

Scalarization converts a vector optimization problem into a scalar
optimization problem.

A scalarizing function

\[
\phi
\]

produces

\[
\boxed{
\min_{\mathbf{x}\in\mathcal{F}}
\phi
(
f_1(\mathbf{x}),
\ldots,
f_k(\mathbf{x})
).
}
\]

Different scalarizations represent different notions of preference and
tradeoff.

%%%%%%%%%%%%%%%%%%%%%%%%%%%%%%%%%%%%%%%%%%%%%%%%%%%%%%%%%%%%
\subsection{Weighted-Sum Method}
\label{subsec:weighted-sum-method}
%%%%%%%%%%%%%%%%%%%%%%%%%%%%%%%%%%%%%%%%%%%%%%%%%%%%%%%%%%%%

Choose nonnegative weights

\[
w_i\geq0,
\]

typically with

\[
\boxed{
\sum_{i=1}^{k}
w_i=1.
}
\]

Then solve

\[
\boxed{
\min_{\mathbf{x}\in\mathcal{F}}
\sum_{i=1}^{k}
w_i
f_i(\mathbf{x}).
}
\]

Larger \(w_i\) places greater emphasis on objective \(i\).

%%%%%%%%%%%%%%%%%%%%%%%%%%%%%%%%%%%%%%%%%%%%%%%%%%%%%%%%%%%%
\subsection{Weighted-Sum Pareto Property}
\label{subsec:weighted-sum-pareto-property}
%%%%%%%%%%%%%%%%%%%%%%%%%%%%%%%%%%%%%%%%%%%%%%%%%%%%%%%%%%%%

If all weights satisfy

\[
\boxed{
w_i>0,
}
\]

then any optimizer of the weighted-sum problem is Pareto optimal under
standard assumptions on feasibility and finite objective values.

If some weights are zero, the solution is generally guaranteed only to be
weakly Pareto optimal without additional conditions.

%%%%%%%%%%%%%%%%%%%%%%%%%%%%%%%%%%%%%%%%%%%%%%%%%%%%%%%%%%%%
\subsection{Worked Example: Weighted Sum}
\label{subsec:worked-weighted-sum}
%%%%%%%%%%%%%%%%%%%%%%%%%%%%%%%%%%%%%%%%%%%%%%%%%%%%%%%%%%%%

\begin{workedexamplebox}{Balancing Cost and Time}

Suppose two feasible plans have objective vectors

\[
A=(100,8),
\]

and

\[
B=(120,5),
\]

where the first component is cost and the second is time.

Use

\[
w_1=0.4,
\qquad
w_2=0.6.
\]

The weighted scores are

\[
S_A
=
0.4(100)
+
0.6(8)
=
44.8,
\]

and

\[
S_B
=
0.4(120)
+
0.6(5)
=
51.
\]

Thus, under these unnormalized weights,

\begin{answerbox}
\[
\boxed{
A
}
\]

has the smaller weighted score.
\end{answerbox}

This example also illustrates why normalization may be important when
objectives use different units.

\end{workedexamplebox}

%%%%%%%%%%%%%%%%%%%%%%%%%%%%%%%%%%%%%%%%%%%%%%%%%%%%%%%%%%%%
\subsection{Supported Pareto Points}
\label{subsec:supported-pareto-points}
%%%%%%%%%%%%%%%%%%%%%%%%%%%%%%%%%%%%%%%%%%%%%%%%%%%%%%%%%%%%

A supported Pareto point is a Pareto-optimal objective vector that solves
a weighted-sum problem for some nonnegative weight vector.

Geometrically, it lies on a supporting hyperplane of the attainable
objective set.

%%%%%%%%%%%%%%%%%%%%%%%%%%%%%%%%%%%%%%%%%%%%%%%%%%%%%%%%%%%%
\subsection{Unsupported Pareto Points}
\label{subsec:unsupported-pareto-points}
%%%%%%%%%%%%%%%%%%%%%%%%%%%%%%%%%%%%%%%%%%%%%%%%%%%%%%%%%%%%

If the attainable objective set is nonconvex, some Pareto-optimal points
may not minimize any weighted sum.

These are unsupported Pareto points.

Thus,

\[
\boxed{
\text{varying weighted sums may fail to generate the entire Pareto front}.
}
\]

%%%%%%%%%%%%%%%%%%%%%%%%%%%%%%%%%%%%%%%%%%%%%%%%%%%%%%%%%%%%
\subsection{Convex Objective Sets}
\label{subsec:convex-objective-sets}
%%%%%%%%%%%%%%%%%%%%%%%%%%%%%%%%%%%%%%%%%%%%%%%%%%%%%%%%%%%%

When the relevant attainable objective region has suitable convexity,
weighted-sum scalarization can recover supported Pareto solutions broadly.

In nonconvex multiobjective problems, additional scalarization techniques
are necessary to reveal unsupported tradeoffs.

%%%%%%%%%%%%%%%%%%%%%%%%%%%%%%%%%%%%%%%%%%%%%%%%%%%%%%%%%%%%
\subsection{Epsilon-Constraint Method}
\label{subsec:epsilon-constraint-method}
%%%%%%%%%%%%%%%%%%%%%%%%%%%%%%%%%%%%%%%%%%%%%%%%%%%%%%%%%%%%

Choose one objective as the primary objective.

Convert the remaining objectives into constraints.

For example,

\[
\boxed{
\begin{aligned}
\min_{\mathbf{x}}
\quad&
f_1(\mathbf{x})
\\
\text{subject to}
\quad&
f_2(\mathbf{x})
\leq
\varepsilon_2,
\\
&
\vdots
\\
&
f_k(\mathbf{x})
\leq
\varepsilon_k,
\\
&
\mathbf{x}\in\mathcal{F}.
\end{aligned}
}
\]

Changing the \(\varepsilon_i\) values generates different tradeoffs.

%%%%%%%%%%%%%%%%%%%%%%%%%%%%%%%%%%%%%%%%%%%%%%%%%%%%%%%%%%%%
\subsection{Advantages of the Epsilon-Constraint Method}
\label{subsec:epsilon-constraint-advantages}
%%%%%%%%%%%%%%%%%%%%%%%%%%%%%%%%%%%%%%%%%%%%%%%%%%%%%%%%%%%%

The epsilon-constraint method can:

\[
\boxed{
\text{generate unsupported Pareto points},
}
\]

\[
\boxed{
\text{preserve the units of the primary objective},
}
\]

and

\[
\boxed{
\text{express explicit acceptable limits on secondary objectives}.
}
\]

It is especially useful when one objective is naturally primary.

%%%%%%%%%%%%%%%%%%%%%%%%%%%%%%%%%%%%%%%%%%%%%%%%%%%%%%%%%%%%
\subsection{Worked Example: Epsilon Constraint}
\label{subsec:worked-epsilon-constraint}
%%%%%%%%%%%%%%%%%%%%%%%%%%%%%%%%%%%%%%%%%%%%%%%%%%%%%%%%%%%%

\begin{workedexamplebox}{Minimum Cost Subject to Emission Limit}

Suppose

\[
C(x)
=
\text{cost},
\]

and

\[
E(x)
=
\text{emissions}.
\]

Instead of combining them with arbitrary weights, solve

\[
\boxed{
\begin{aligned}
\min_x
\quad&
C(x)
\\
\text{subject to}
\quad&
E(x)\leq\varepsilon,
\\
&
x\in\mathcal{F}.
\end{aligned}
}
\]

Reducing \(\varepsilon\) imposes a stricter emissions requirement.

The resulting sequence of solutions traces the cost--emissions tradeoff.

\end{workedexamplebox}

%%%%%%%%%%%%%%%%%%%%%%%%%%%%%%%%%%%%%%%%%%%%%%%%%%%%%%%%%%%%
\subsection{Lexicographic Optimization}
\label{subsec:lexicographic-optimization}
%%%%%%%%%%%%%%%%%%%%%%%%%%%%%%%%%%%%%%%%%%%%%%%%%%%%%%%%%%%%

Lexicographic optimization ranks objectives by priority.

First optimize

\[
f_1.
\]

Among all solutions optimal for \(f_1\), optimize

\[
f_2.
\]

Then optimize \(f_3\), and so on.

Thus

\[
\boxed{
f_1
\succ
f_2
\succ
\cdots
\succ
f_k.
}
\]

%%%%%%%%%%%%%%%%%%%%%%%%%%%%%%%%%%%%%%%%%%%%%%%%%%%%%%%%%%%%
\subsection{Lexicographic Formulation}
\label{subsec:lexicographic-formulation}
%%%%%%%%%%%%%%%%%%%%%%%%%%%%%%%%%%%%%%%%%%%%%%%%%%%%%%%%%%%%

Suppose

\[
f_1^*
=
\min_{\mathbf{x}\in\mathcal{F}}
f_1(\mathbf{x}).
\]

The second problem is

\[
\boxed{
\begin{aligned}
\min_{\mathbf{x}}
\quad&
f_2(\mathbf{x})
\\
\text{subject to}
\quad&
f_1(\mathbf{x})
=
f_1^*,
\\
&
\mathbf{x}\in\mathcal{F}.
\end{aligned}
}
\]

This process continues through the ordered objectives.

%%%%%%%%%%%%%%%%%%%%%%%%%%%%%%%%%%%%%%%%%%%%%%%%%%%%%%%%%%%%
\subsection{Priority Tolerances}
\label{subsec:lexicographic-tolerances}
%%%%%%%%%%%%%%%%%%%%%%%%%%%%%%%%%%%%%%%%%%%%%%%%%%%%%%%%%%%%

In practice, requiring exact preservation of a higher-priority objective
may be unnecessarily rigid.

One may allow

\[
\boxed{
f_1(\mathbf{x})
\leq
f_1^*+\delta_1,
}
\]

where \(\delta_1\) is an acceptable degradation.

This creates controlled tradeoffs while retaining priority structure.

%%%%%%%%%%%%%%%%%%%%%%%%%%%%%%%%%%%%%%%%%%%%%%%%%%%%%%%%%%%%
\subsection{Goal Programming}
\label{subsec:goal-programming}
%%%%%%%%%%%%%%%%%%%%%%%%%%%%%%%%%%%%%%%%%%%%%%%%%%%%%%%%%%%%

Goal programming specifies desired target values for objectives.

Suppose objective \(i\) has target

\[
g_i.
\]

Introduce deviations

\[
d_i^-\geq0,
\qquad
d_i^+\geq0,
\]

such that

\[
\boxed{
f_i(\mathbf{x})
+
d_i^-
-
d_i^+
=
g_i.
}
\]

%%%%%%%%%%%%%%%%%%%%%%%%%%%%%%%%%%%%%%%%%%%%%%%%%%%%%%%%%%%%
\subsection{Meaning of Goal Deviations}
\label{subsec:goal-deviations}
%%%%%%%%%%%%%%%%%%%%%%%%%%%%%%%%%%%%%%%%%%%%%%%%%%%%%%%%%%%%

The variable

\[
d_i^-
\]

measures underachievement relative to the target.

The variable

\[
d_i^+
\]

measures overachievement.

Depending on the objective, one or both deviations may be undesirable.

%%%%%%%%%%%%%%%%%%%%%%%%%%%%%%%%%%%%%%%%%%%%%%%%%%%%%%%%%%%%
\subsection{Weighted Goal Programming}
\label{subsec:weighted-goal-programming}
%%%%%%%%%%%%%%%%%%%%%%%%%%%%%%%%%%%%%%%%%%%%%%%%%%%%%%%%%%%%

A weighted goal-programming objective may be

\[
\boxed{
\min
\sum_{i=1}^{k}
\left(
w_i^-d_i^-
+
w_i^+d_i^+
\right).
}
\]

The weights specify the relative penalty associated with missing different
targets.

%%%%%%%%%%%%%%%%%%%%%%%%%%%%%%%%%%%%%%%%%%%%%%%%%%%%%%%%%%%%
\subsection{Worked Example: Goal Programming}
\label{subsec:worked-goal-programming}
%%%%%%%%%%%%%%%%%%%%%%%%%%%%%%%%%%%%%%%%%%%%%%%%%%%%%%%%%%%%

\begin{workedexamplebox}{Production Targets}

Suppose desired profit is at least

\[
100,
\]

while desired overtime is at most

\[
20.
\]

Introduce deviations satisfying

\[
P(x)
+
d_1^-
-
d_1^+
=
100,
\]

and

\[
O(x)
+
d_2^-
-
d_2^+
=
20.
\]

If missing the profit target and exceeding overtime are undesirable, an
objective can penalize

\[
\boxed{
d_1^-
}
\]

and

\[
\boxed{
d_2^+.
}
\]

For example,

\[
\boxed{
\min
5d_1^-
+
2d_2^+.
}
\]

\end{workedexamplebox}

%%%%%%%%%%%%%%%%%%%%%%%%%%%%%%%%%%%%%%%%%%%%%%%%%%%%%%%%%%%%
\subsection{Preemptive Goal Programming}
\label{subsec:preemptive-goal-programming}
%%%%%%%%%%%%%%%%%%%%%%%%%%%%%%%%%%%%%%%%%%%%%%%%%%%%%%%%%%%%

Preemptive goal programming assigns priority levels.

A higher-priority goal is optimized before any lower-priority goal is
considered.

This resembles lexicographic optimization.

It is useful when objectives are not merely weighted differently but have
strict institutional priorities.

%%%%%%%%%%%%%%%%%%%%%%%%%%%%%%%%%%%%%%%%%%%%%%%%%%%%%%%%%%%%
\subsection{Reference-Point Methods}
\label{subsec:reference-point-methods}
%%%%%%%%%%%%%%%%%%%%%%%%%%%%%%%%%%%%%%%%%%%%%%%%%%%%%%%%%%%%

A decision maker may specify a desired objective vector

\[
\boxed{
\mathbf{r}
=
(r_1,\ldots,r_k).
}
\]

Optimization then seeks a feasible objective vector close to
\(\mathbf{r}\).

A simple model is

\[
\boxed{
\min_{\mathbf{x}\in\mathcal{F}}
\|
\mathbf{f}(\mathbf{x})
-
\mathbf{r}
\|.
}
\]

%%%%%%%%%%%%%%%%%%%%%%%%%%%%%%%%%%%%%%%%%%%%%%%%%%%%%%%%%%%%
\subsection{Reference Points Need Not Be Feasible}
\label{subsec:reference-points-feasibility}
%%%%%%%%%%%%%%%%%%%%%%%%%%%%%%%%%%%%%%%%%%%%%%%%%%%%%%%%%%%%

The reference vector may represent aspirations.

It need not be attainable.

Optimization identifies a feasible compromise near the desired target.

This makes reference-point approaches useful in interactive decision
support.

%%%%%%%%%%%%%%%%%%%%%%%%%%%%%%%%%%%%%%%%%%%%%%%%%%%%%%%%%%%%
\subsection{Compromise Programming}
\label{subsec:compromise-programming}
%%%%%%%%%%%%%%%%%%%%%%%%%%%%%%%%%%%%%%%%%%%%%%%%%%%%%%%%%%%%

Compromise programming minimizes distance from the ideal objective vector.

A weighted \(p\)-norm formulation is

\[
\boxed{
\min_{\mathbf{x}\in\mathcal{F}}
\left[
\sum_{i=1}^{k}
w_i
\left|
\frac{
f_i(\mathbf{x})
-
z_i^{\mathrm{ideal}}
}{
s_i
}
\right|^p
\right]^{1/p},
}
\]

where \(s_i\) are scaling factors.

%%%%%%%%%%%%%%%%%%%%%%%%%%%%%%%%%%%%%%%%%%%%%%%%%%%%%%%%%%%%
\subsection{Chebyshev Scalarization}
\label{subsec:chebyshev-scalarization}
%%%%%%%%%%%%%%%%%%%%%%%%%%%%%%%%%%%%%%%%%%%%%%%%%%%%%%%%%%%%

Using the infinity norm gives a weighted Chebyshev scalarization:

\[
\boxed{
\min_{\mathbf{x}\in\mathcal{F}}
\max_i
\left\{
w_i
\left|
f_i(\mathbf{x})
-
z_i^{\mathrm{ref}}
\right|
\right\}.
}
\]

This minimizes the largest weighted deviation from the reference point.

%%%%%%%%%%%%%%%%%%%%%%%%%%%%%%%%%%%%%%%%%%%%%%%%%%%%%%%%%%%%
\subsection{Epigraph Form of Chebyshev Scalarization}
\label{subsec:chebyshev-epigraph}
%%%%%%%%%%%%%%%%%%%%%%%%%%%%%%%%%%%%%%%%%%%%%%%%%%%%%%%%%%%%

Introduce variable \(t\).

Then

\[
\boxed{
\begin{aligned}
\min_{\mathbf{x},t}
\quad&
t
\\
\text{subject to}
\quad&
w_i
\left[
f_i(\mathbf{x})
-
z_i^{\mathrm{ref}}
\right]
\leq
t,
\qquad
i=1,\ldots,k,
\\
&
\mathbf{x}\in\mathcal{F},
\end{aligned}
}
\]

for an appropriate one-sided minimization formulation.

This converts the maximum into explicit constraints.

%%%%%%%%%%%%%%%%%%%%%%%%%%%%%%%%%%%%%%%%%%%%%%%%%%%%%%%%%%%%
\subsection{Achievement Scalarizing Functions}
\label{subsec:achievement-scalarizing-functions}
%%%%%%%%%%%%%%%%%%%%%%%%%%%%%%%%%%%%%%%%%%%%%%%%%%%%%%%%%%%%

An achievement scalarizing function measures how well a feasible solution
achieves aspiration levels.

A common structure combines:

\[
\boxed{
\text{maximum weighted deviation}
}
\]

with

\[
\boxed{
\text{a small augmentation term}.
}
\]

The augmentation helps avoid weakly Pareto-optimal solutions.

%%%%%%%%%%%%%%%%%%%%%%%%%%%%%%%%%%%%%%%%%%%%%%%%%%%%%%%%%%%%
\subsection{Multiobjective Linear Programming}
\label{subsec:multiobjective-linear-programming}
%%%%%%%%%%%%%%%%%%%%%%%%%%%%%%%%%%%%%%%%%%%%%%%%%%%%%%%%%%%%

A multiobjective linear program has linear objectives and linear
constraints:

\[
\boxed{
\min
\begin{pmatrix}
c_1^Tx\\
c_2^Tx\\
\vdots\\
c_k^Tx
\end{pmatrix}
}
\]

subject to

\[
\boxed{
Ax\leq b.
}
\]

Weighted sums produce ordinary linear programs.

Epsilon-constraint formulations also remain linear.

%%%%%%%%%%%%%%%%%%%%%%%%%%%%%%%%%%%%%%%%%%%%%%%%%%%%%%%%%%%%
\subsection{Two-Objective Linear Example}
\label{subsec:two-objective-linear-example}
%%%%%%%%%%%%%%%%%%%%%%%%%%%%%%%%%%%%%%%%%%%%%%%%%%%%%%%%%%%%

Suppose

\[
\boxed{
\min
\begin{pmatrix}
2x+y\\
x+3y
\end{pmatrix}
}
\]

subject to

\[
x+y\geq4,
\]

\[
x,y\geq0.
\]

A weighted sum with

\[
w\in[0,1]
\]

becomes

\[
\boxed{
\min
\left[
w(2x+y)
+
(1-w)(x+3y)
\right].
}
\]

Different \(w\) values produce different efficient solutions.

%%%%%%%%%%%%%%%%%%%%%%%%%%%%%%%%%%%%%%%%%%%%%%%%%%%%%%%%%%%%
\subsection{Multiobjective Convex Optimization}
\label{subsec:multiobjective-convex-optimization}
%%%%%%%%%%%%%%%%%%%%%%%%%%%%%%%%%%%%%%%%%%%%%%%%%%%%%%%%%%%%

Suppose:

\[
\mathcal{F}
\]

is convex and each

\[
f_i
\]

is convex.

Then the problem is a convex multiobjective optimization problem.

Weighted sums with nonnegative weights produce convex scalar optimization
problems.

This gives strong computational advantages.

%%%%%%%%%%%%%%%%%%%%%%%%%%%%%%%%%%%%%%%%%%%%%%%%%%%%%%%%%%%%
\subsection{Pareto Optimality in Convex Problems}
\label{subsec:pareto-convex-problems}
%%%%%%%%%%%%%%%%%%%%%%%%%%%%%%%%%%%%%%%%%%%%%%%%%%%%%%%%%%%%

Under suitable convexity assumptions, Pareto-optimal points can often be
characterized using positive weighted combinations of objective
gradients.

This connects multiobjective optimization with KKT theory.

%%%%%%%%%%%%%%%%%%%%%%%%%%%%%%%%%%%%%%%%%%%%%%%%%%%%%%%%%%%%
\subsection{First-Order Pareto Condition}
\label{subsec:first-order-pareto-condition}
%%%%%%%%%%%%%%%%%%%%%%%%%%%%%%%%%%%%%%%%%%%%%%%%%%%%%%%%%%%%

For an unconstrained differentiable problem, a necessary condition for a
local Pareto optimum under suitable regularity is existence of weights

\[
\lambda_i\geq0,
\]

not all zero, such that

\[
\boxed{
\sum_{i=1}^{k}
\lambda_i
\nabla f_i(\mathbf{x}^*)
=
0.
}
\]

The condition says there is no common first-order direction improving all
objectives.

%%%%%%%%%%%%%%%%%%%%%%%%%%%%%%%%%%%%%%%%%%%%%%%%%%%%%%%%%%%%
\subsection{Normalized Objective Weights}
\label{subsec:normalized-objective-weights}
%%%%%%%%%%%%%%%%%%%%%%%%%%%%%%%%%%%%%%%%%%%%%%%%%%%%%%%%%%%%

The weights may be normalized so that

\[
\boxed{
\lambda_i\geq0,
\qquad
\sum_i\lambda_i=1.
}
\]

Then

\[
\boxed{
\sum_i
\lambda_i
\nabla f_i(\mathbf{x}^*)
=
0.
}
\]

This resembles stationarity of a weighted scalar objective.

%%%%%%%%%%%%%%%%%%%%%%%%%%%%%%%%%%%%%%%%%%%%%%%%%%%%%%%%%%%%
\subsection{Constrained Multiobjective KKT Condition}
\label{subsec:multiobjective-kkt}
%%%%%%%%%%%%%%%%%%%%%%%%%%%%%%%%%%%%%%%%%%%%%%%%%%%%%%%%%%%%

Consider constraints

\[
g_j(x)\leq0,
\]

and

\[
h_\ell(x)=0.
\]

A multiobjective first-order condition can take form

\[
\boxed{
\sum_i
\lambda_i
\nabla f_i(x^*)
+
\sum_j
\mu_j
\nabla g_j(x^*)
+
\sum_\ell
\nu_\ell
\nabla h_\ell(x^*)
=
0,
}
\]

with

\[
\boxed{
\lambda_i\geq0,
}
\]

\[
\boxed{
\mu_j\geq0,
}
\]

and complementary slackness

\[
\boxed{
\mu_jg_j(x^*)=0.
}
\]

Suitable regularity conditions are required.

%%%%%%%%%%%%%%%%%%%%%%%%%%%%%%%%%%%%%%%%%%%%%%%%%%%%%%%%%%%%
\subsection{Marginal Tradeoff Interpretation}
\label{subsec:marginal-tradeoff}
%%%%%%%%%%%%%%%%%%%%%%%%%%%%%%%%%%%%%%%%%%%%%%%%%%%%%%%%%%%%

For two smooth objectives along a differentiable Pareto frontier, one may
study the local tradeoff

\[
\boxed{
\frac{
df_2
}{
df_1
}.
}
\]

This measures how much one objective changes per unit change in another.

Such rates can support decision-maker interpretation.

%%%%%%%%%%%%%%%%%%%%%%%%%%%%%%%%%%%%%%%%%%%%%%%%%%%%%%%%%%%%
\subsection{Proper Efficiency}
\label{subsec:proper-efficiency}
%%%%%%%%%%%%%%%%%%%%%%%%%%%%%%%%%%%%%%%%%%%%%%%%%%%%%%%%%%%%

Some Pareto points involve extreme tradeoffs.

For example, an arbitrarily small improvement in one objective may require
an arbitrarily large deterioration in another.

Proper efficiency concepts exclude certain unbounded tradeoff behavior.

Several formal definitions exist, including Geoffrion proper efficiency.

%%%%%%%%%%%%%%%%%%%%%%%%%%%%%%%%%%%%%%%%%%%%%%%%%%%%%%%%%%%%
\subsection{Preference Information}
\label{subsec:preference-information}
%%%%%%%%%%%%%%%%%%%%%%%%%%%%%%%%%%%%%%%%%%%%%%%%%%%%%%%%%%%%

Mathematics can identify feasible tradeoffs.

It does not determine automatically which tradeoff a decision maker should
prefer.

Preference information may include:

\[
\boxed{
\text{weights},
}
\]

\[
\boxed{
\text{priorities},
}
\]

\[
\boxed{
\text{targets},
}
\]

\[
\boxed{
\text{acceptable thresholds},
}
\]

or

\[
\boxed{
\text{utility functions}.
}
\]

%%%%%%%%%%%%%%%%%%%%%%%%%%%%%%%%%%%%%%%%%%%%%%%%%%%%%%%%%%%%
\subsection{A Priori Methods}
\label{subsec:a-priori-multiobjective}
%%%%%%%%%%%%%%%%%%%%%%%%%%%%%%%%%%%%%%%%%%%%%%%%%%%%%%%%%%%%

In an a priori method, preferences are specified before optimization.

Examples include:

\[
\boxed{
\text{weighted sums},
}
\]

\[
\boxed{
\text{lexicographic priorities},
}
\]

and

\[
\boxed{
\text{goal-programming weights}.
}
\]

The optimization then produces one preferred compromise.

%%%%%%%%%%%%%%%%%%%%%%%%%%%%%%%%%%%%%%%%%%%%%%%%%%%%%%%%%%%%
\subsection{A Posteriori Methods}
\label{subsec:a-posteriori-multiobjective}
%%%%%%%%%%%%%%%%%%%%%%%%%%%%%%%%%%%%%%%%%%%%%%%%%%%%%%%%%%%%

In an a posteriori method, optimization first generates many Pareto
solutions.

The decision maker then examines the tradeoff set and selects a preferred
solution.

This approach is useful when preferences are difficult to specify in
advance.

%%%%%%%%%%%%%%%%%%%%%%%%%%%%%%%%%%%%%%%%%%%%%%%%%%%%%%%%%%%%
\subsection{Interactive Methods}
\label{subsec:interactive-multiobjective}
%%%%%%%%%%%%%%%%%%%%%%%%%%%%%%%%%%%%%%%%%%%%%%%%%%%%%%%%%%%%

Interactive methods alternate between:

\[
\boxed{
\text{optimization}
}
\]

and

\[
\boxed{
\text{preference feedback}.
}
\]

The decision maker may adjust aspirations, weights, or acceptable
tradeoffs after seeing intermediate solutions.

This can reduce the number of Pareto points that must be generated.

%%%%%%%%%%%%%%%%%%%%%%%%%%%%%%%%%%%%%%%%%%%%%%%%%%%%%%%%%%%%
\subsection{Generating a Pareto Frontier}
\label{subsec:generating-pareto-frontier}
%%%%%%%%%%%%%%%%%%%%%%%%%%%%%%%%%%%%%%%%%%%%%%%%%%%%%%%%%%%%

For two objectives, a common strategy is:

\begin{enumerate}
    \item solve each objective separately;
    \item determine the objective ranges;
    \item select a sequence of weights or epsilon values;
    \item solve the corresponding scalar problems;
    \item discard dominated or duplicate points;
    \item order the remaining solutions along the Pareto frontier.
\end{enumerate}

%%%%%%%%%%%%%%%%%%%%%%%%%%%%%%%%%%%%%%%%%%%%%%%%%%%%%%%%%%%%
\subsection{Pareto Frontier Resolution}
\label{subsec:pareto-frontier-resolution}
%%%%%%%%%%%%%%%%%%%%%%%%%%%%%%%%%%%%%%%%%%%%%%%%%%%%%%%%%%%%

A coarse set of scalarization parameters may miss important regions of the
frontier.

A very fine grid may require many optimization solves.

Adaptive procedures focus additional computations where the frontier
changes rapidly.

%%%%%%%%%%%%%%%%%%%%%%%%%%%%%%%%%%%%%%%%%%%%%%%%%%%%%%%%%%%%
\subsection{Knee Points}
\label{subsec:knee-points}
%%%%%%%%%%%%%%%%%%%%%%%%%%%%%%%%%%%%%%%%%%%%%%%%%%%%%%%%%%%%

A knee point is an informal term for a Pareto solution where a small
improvement in one objective requires a large deterioration in another.

Such points can be attractive compromise solutions.

However, the precise definition of a knee depends on scaling and the
chosen geometry.

%%%%%%%%%%%%%%%%%%%%%%%%%%%%%%%%%%%%%%%%%%%%%%%%%%%%%%%%%%%%
\subsection{Discrete Multiobjective Optimization}
\label{subsec:discrete-multiobjective-optimization}
%%%%%%%%%%%%%%%%%%%%%%%%%%%%%%%%%%%%%%%%%%%%%%%%%%%%%%%%%%%%

If decisions are integer or combinatorial, the Pareto frontier may consist
of isolated objective vectors.

For example,

\[
\boxed{
\min
\begin{pmatrix}
c^Tx\\
d^Tx
\end{pmatrix}
}
\]

subject to

\[
\boxed{
Ax\leq b,
\qquad
x\in\{0,1\}^n
}
\]

is a multiobjective binary optimization problem.

%%%%%%%%%%%%%%%%%%%%%%%%%%%%%%%%%%%%%%%%%%%%%%%%%%%%%%%%%%%%
\subsection{Nondominated Set Enumeration}
\label{subsec:nondominated-set-enumeration}
%%%%%%%%%%%%%%%%%%%%%%%%%%%%%%%%%%%%%%%%%%%%%%%%%%%%%%%%%%%%

For discrete problems, one may seek every nondominated objective vector.

The number of such points can itself be very large.

Therefore full enumeration may be computationally expensive even when
each scalar subproblem is manageable.

%%%%%%%%%%%%%%%%%%%%%%%%%%%%%%%%%%%%%%%%%%%%%%%%%%%%%%%%%%%%
\subsection{Evolutionary Multiobjective Optimization}
\label{subsec:evolutionary-multiobjective-optimization}
%%%%%%%%%%%%%%%%%%%%%%%%%%%%%%%%%%%%%%%%%%%%%%%%%%%%%%%%%%%%

Evolutionary algorithms maintain a population of candidate solutions.

Instead of converging toward one objective value, they attempt to maintain
a diverse set of nondominated solutions.

Important algorithm families include methods based on:

\[
\boxed{
\text{nondominated sorting},
}
\]

\[
\boxed{
\text{crowding measures},
}
\]

and

\[
\boxed{
\text{reference directions}.
}
\]

%%%%%%%%%%%%%%%%%%%%%%%%%%%%%%%%%%%%%%%%%%%%%%%%%%%%%%%%%%%%
\subsection{Nondominated Sorting}
\label{subsec:nondominated-sorting}
%%%%%%%%%%%%%%%%%%%%%%%%%%%%%%%%%%%%%%%%%%%%%%%%%%%%%%%%%%%%

Given a population of objective vectors, the first nondominated front
contains all points not dominated by any other point.

Remove those points.

The next nondominated set forms the second front.

Repeating produces dominance levels.

%%%%%%%%%%%%%%%%%%%%%%%%%%%%%%%%%%%%%%%%%%%%%%%%%%%%%%%%%%%%
\subsection{Diversity}
\label{subsec:pareto-diversity}
%%%%%%%%%%%%%%%%%%%%%%%%%%%%%%%%%%%%%%%%%%%%%%%%%%%%%%%%%%%%

An approximate Pareto set should ideally have both:

\[
\boxed{
\text{good convergence toward the true frontier},
}
\]

and

\[
\boxed{
\text{good coverage across the frontier}.
}
\]

A cluster of nearly identical solutions may approximate only a small
portion of the tradeoff space.

%%%%%%%%%%%%%%%%%%%%%%%%%%%%%%%%%%%%%%%%%%%%%%%%%%%%%%%%%%%%
\subsection{Hypervolume Indicator}
\label{subsec:hypervolume-indicator}
%%%%%%%%%%%%%%%%%%%%%%%%%%%%%%%%%%%%%%%%%%%%%%%%%%%%%%%%%%%%

The hypervolume indicator measures the region of objective space dominated
by a set of nondominated solutions relative to a chosen reference point.

Larger hypervolume generally indicates some combination of:

\[
\boxed{
\text{better objective quality}
}
\]

and

\[
\boxed{
\text{broader Pareto-front coverage}.
}
\]

Its value depends on the reference point and scaling.

%%%%%%%%%%%%%%%%%%%%%%%%%%%%%%%%%%%%%%%%%%%%%%%%%%%%%%%%%%%%
\subsection{Weighted Sum Sensitivity}
\label{subsec:weight-sensitivity}
%%%%%%%%%%%%%%%%%%%%%%%%%%%%%%%%%%%%%%%%%%%%%%%%%%%%%%%%%%%%

A small change in weights can sometimes produce a large change in the
chosen solution.

This is especially common when:

\[
\boxed{
\text{the Pareto front is sharply curved},
}
\]

or

\[
\boxed{
\text{the feasible decisions are discrete}.
}
\]

Weight sensitivity should therefore be examined rather than reporting a
single arbitrary weight vector.

%%%%%%%%%%%%%%%%%%%%%%%%%%%%%%%%%%%%%%%%%%%%%%%%%%%%%%%%%%%%
\subsection{Objective Scaling and Units}
\label{subsec:objective-scaling-units}
%%%%%%%%%%%%%%%%%%%%%%%%%%%%%%%%%%%%%%%%%%%%%%%%%%%%%%%%%%%%

Suppose cost is measured in dollars and time in hours.

Then the expression

\[
\boxed{
0.5
(\text{cost})
+
0.5
(\text{time})
}
\]

does not automatically imply equal importance.

The numerical scales and units differ.

Normalization or explicit economic conversion may be required.

%%%%%%%%%%%%%%%%%%%%%%%%%%%%%%%%%%%%%%%%%%%%%%%%%%%%%%%%%%%%
\subsection{Fairness as a Multiobjective Problem}
\label{subsec:fairness-multiobjective}
%%%%%%%%%%%%%%%%%%%%%%%%%%%%%%%%%%%%%%%%%%%%%%%%%%%%%%%%%%%%

Some allocation problems must balance:

\[
\boxed{
\text{efficiency}
}
\]

against

\[
\boxed{
\text{fairness}.
}
\]

Alternatively, different groups may have separate performance objectives.

A multiobjective model makes these competing criteria explicit instead of
hiding them inside one scalar score.

%%%%%%%%%%%%%%%%%%%%%%%%%%%%%%%%%%%%%%%%%%%%%%%%%%%%%%%%%%%%
\subsection{Max--Min Fairness Preview}
\label{subsec:max-min-fairness}
%%%%%%%%%%%%%%%%%%%%%%%%%%%%%%%%%%%%%%%%%%%%%%%%%%%%%%%%%%%%

Suppose each stakeholder receives utility

\[
u_i(x).
\]

A max--min fairness objective is

\[
\boxed{
\max_x
\min_i
u_i(x).
}
\]

Introduce \(t\) and write

\[
\boxed{
\begin{aligned}
\max_{x,t}
\quad&
t
\\
\text{subject to}
\quad&
u_i(x)\geq t,
\qquad
i=1,\ldots,k.
\end{aligned}
}
\]

This emphasizes the worst-off stakeholder.

%%%%%%%%%%%%%%%%%%%%%%%%%%%%%%%%%%%%%%%%%%%%%%%%%%%%%%%%%%%%
\subsection{Sustainability Optimization}
\label{subsec:sustainability-multiobjective}
%%%%%%%%%%%%%%%%%%%%%%%%%%%%%%%%%%%%%%%%%%%%%%%%%%%%%%%%%%%%

Sustainability planning often balances:

\[
\boxed{
\text{economic cost},
}
\]

\[
\boxed{
\text{environmental impact},
}
\]

and

\[
\boxed{
\text{social outcomes}.
}
\]

A multiobjective model can expose the tradeoffs among these criteria rather
than requiring an arbitrary single conversion factor.

%%%%%%%%%%%%%%%%%%%%%%%%%%%%%%%%%%%%%%%%%%%%%%%%%%%%%%%%%%%%
\subsection{Portfolio Tradeoff}
\label{subsec:portfolio-multiobjective}
%%%%%%%%%%%%%%%%%%%%%%%%%%%%%%%%%%%%%%%%%%%%%%%%%%%%%%%%%%%%

Portfolio selection naturally contains competing objectives:

\[
\boxed{
\text{maximize expected return},
}
\]

and

\[
\boxed{
\text{minimize risk}.
}
\]

One formulation is

\[
\boxed{
\min_w
\begin{pmatrix}
-w^T\mu\\
w^T\Sigma w
\end{pmatrix}
}
\]

subject to

\[
\boxed{
\mathbf{1}^Tw=1.
}
\]

The resulting Pareto frontier corresponds to an efficient risk--return
frontier.

%%%%%%%%%%%%%%%%%%%%%%%%%%%%%%%%%%%%%%%%%%%%%%%%%%%%%%%%%%%%
\subsection{Engineering Design Tradeoffs}
\label{subsec:engineering-multiobjective}
%%%%%%%%%%%%%%%%%%%%%%%%%%%%%%%%%%%%%%%%%%%%%%%%%%%%%%%%%%%%

An engineering design may simultaneously seek to minimize:

\[
\boxed{
\text{mass},
}
\]

\[
\boxed{
\text{cost},
}
\]

and

\[
\boxed{
\text{energy consumption},
}
\]

while maximizing:

\[
\boxed{
\text{reliability}
}
\]

or

\[
\boxed{
\text{performance}.
}
\]

Multiobjective optimization helps identify designs where no criterion can
be improved without sacrificing another.

%%%%%%%%%%%%%%%%%%%%%%%%%%%%%%%%%%%%%%%%%%%%%%%%%%%%%%%%%%%%
\subsection{Operations Research Applications}
\label{subsec:or-multiobjective}
%%%%%%%%%%%%%%%%%%%%%%%%%%%%%%%%%%%%%%%%%%%%%%%%%%%%%%%%%%%%

Operations-research models may balance:

\[
\boxed{
\text{service quality},
}
\]

\[
\boxed{
\text{operating cost},
}
\]

\[
\boxed{
\text{risk},
}
\]

\[
\boxed{
\text{travel distance},
}
\]

and

\[
\boxed{
\text{resource utilization}.
}
\]

These objectives frequently represent different stakeholders.

%%%%%%%%%%%%%%%%%%%%%%%%%%%%%%%%%%%%%%%%%%%%%%%%%%%%%%%%%%%%
\subsection{Multiobjective Optimization Workflow}
\label{subsec:multiobjective-workflow}
%%%%%%%%%%%%%%%%%%%%%%%%%%%%%%%%%%%%%%%%%%%%%%%%%%%%%%%%%%%%

A practical workflow is:

\begin{enumerate}
    \item define the decision variables;
    \item define the feasible set;
    \item identify all meaningful objectives;
    \item determine which objectives are minimized and which are
          maximized;
    \item convert directions consistently if needed;
    \item solve each objective individually;
    \item estimate ideal and objective-range information;
    \item normalize objectives when scales differ substantially;
    \item choose a scalarization or frontier-generation method;
    \item generate candidate Pareto solutions;
    \item remove dominated solutions;
    \item visualize objective tradeoffs when dimension permits;
    \item perform sensitivity analysis on weights or thresholds;
    \item obtain preference information from the decision maker;
    \item select a final compromise;
    \item report both the selected solution and the alternatives sacrificed.
\end{enumerate}

%%%%%%%%%%%%%%%%%%%%%%%%%%%%%%%%%%%%%%%%%%%%%%%%%%%%%%%%%%%%
\subsection{Choosing a Multiobjective Method}
\label{subsec:choosing-multiobjective-method}
%%%%%%%%%%%%%%%%%%%%%%%%%%%%%%%%%%%%%%%%%%%%%%%%%%%%%%%%%%%%

When clear numerical tradeoff weights exist, consider:

\[
\boxed{
\text{weighted-sum scalarization}.
}
\]

When one objective is primary and others have limits, consider:

\[
\boxed{
\text{epsilon-constraint optimization}.
}
\]

When goals have explicit targets, consider:

\[
\boxed{
\text{goal programming}.
}
\]

When objectives have strict priorities, consider:

\[
\boxed{
\text{lexicographic optimization}.
}
\]

When preferences are not known initially, consider:

\[
\boxed{
\text{Pareto-front generation}
}
\]

or an interactive method.

%%%%%%%%%%%%%%%%%%%%%%%%%%%%%%%%%%%%%%%%%%%%%%%%%%%%%%%%%%%%
\subsection{Common Errors}
\label{subsec:multiobjective-common-errors}
%%%%%%%%%%%%%%%%%%%%%%%%%%%%%%%%%%%%%%%%%%%%%%%%%%%%%%%%%%%%

\subsubsection{Looking for One Solution That Minimizes Every Objective}

Conflicting objectives generally do not possess a common individual
minimizer.

The goal is usually to find Pareto-efficient tradeoffs.

%%%%%%%%%%%%%%%%%%%%%%%%%%%%%%%%%%%%%%%%%%%%%%%%%%%%%%%%%%%%
\subsubsection{Confusing Pareto Optimality with a Unique Best Solution}

Many Pareto-optimal solutions can exist.

Pareto optimality removes dominated alternatives but does not by itself
select one final solution.

%%%%%%%%%%%%%%%%%%%%%%%%%%%%%%%%%%%%%%%%%%%%%%%%%%%%%%%%%%%%
\subsubsection{Ignoring Optimization Direction}

Dominance definitions depend on whether objectives are minimized or
maximized.

A convenient approach is to convert all objectives to a common direction
before comparing them.

%%%%%%%%%%%%%%%%%%%%%%%%%%%%%%%%%%%%%%%%%%%%%%%%%%%%%%%%%%%%
\subsubsection{Using Weights Without Normalization}

If objectives differ greatly in magnitude, their numerical scales may
dominate the weighted sum regardless of intended preference.

Normalization should be considered.

%%%%%%%%%%%%%%%%%%%%%%%%%%%%%%%%%%%%%%%%%%%%%%%%%%%%%%%%%%%%
\subsubsection{Treating Weights as Percent Importance Automatically}

Weights in

\[
\sum_iw_if_i(x)
\]

interact with units and objective scales.

A weight of \(0.5\) does not necessarily mean an objective has exactly
\(50\%\) practical importance.

%%%%%%%%%%%%%%%%%%%%%%%%%%%%%%%%%%%%%%%%%%%%%%%%%%%%%%%%%%%%
\subsubsection{Assuming Weighted Sums Find Every Pareto Point}

For nonconvex objective sets, weighted sums can miss unsupported Pareto
points.

Use methods such as epsilon constraints when full frontier coverage is
important.

%%%%%%%%%%%%%%%%%%%%%%%%%%%%%%%%%%%%%%%%%%%%%%%%%%%%%%%%%%%%
\subsubsection{Confusing Weak Pareto Optimality with Pareto Optimality}

Weak Pareto optimality only prevents simultaneous strict improvement in
every objective.

Pareto optimality is stronger.

%%%%%%%%%%%%%%%%%%%%%%%%%%%%%%%%%%%%%%%%%%%%%%%%%%%%%%%%%%%%
\subsubsection{Assuming the Ideal Point Is Feasible}

The ideal vector combines separate single-objective optima.

Those values may not be attainable simultaneously.

%%%%%%%%%%%%%%%%%%%%%%%%%%%%%%%%%%%%%%%%%%%%%%%%%%%%%%%%%%%%
\subsubsection{Using an Inaccurate Nadir Point for Scaling}

Nadir estimates can influence normalized tradeoffs significantly.

Poor estimates may distort the apparent geometry of the frontier.

%%%%%%%%%%%%%%%%%%%%%%%%%%%%%%%%%%%%%%%%%%%%%%%%%%%%%%%%%%%%
\subsubsection{Ignoring Duplicate Objective Vectors}

Different decisions can produce the same objective vector.

A Pareto frontier in objective space need not uniquely identify the
decision.

%%%%%%%%%%%%%%%%%%%%%%%%%%%%%%%%%%%%%%%%%%%%%%%%%%%%%%%%%%%%
\subsubsection{Assuming More Pareto Points Always Means a Better Approximation}

A large number of points concentrated in one region may provide poor
coverage.

Both convergence and diversity matter.

%%%%%%%%%%%%%%%%%%%%%%%%%%%%%%%%%%%%%%%%%%%%%%%%%%%%%%%%%%%%
\subsubsection{Using Hypervolume Without Stating the Reference Point}

Hypervolume depends on the chosen reference point.

Results are not meaningful without specifying that reference.

%%%%%%%%%%%%%%%%%%%%%%%%%%%%%%%%%%%%%%%%%%%%%%%%%%%%%%%%%%%%
\subsubsection{Choosing a Final Solution Without Preference Information}

Optimization can reveal efficient tradeoffs.

Selecting among them generally requires values, priorities, or preferences
from the relevant decision maker.

%%%%%%%%%%%%%%%%%%%%%%%%%%%%%%%%%%%%%%%%%%%%%%%%%%%%%%%%%%%%
\subsection{Applications}
\label{subsec:multiobjective-applications}
%%%%%%%%%%%%%%%%%%%%%%%%%%%%%%%%%%%%%%%%%%%%%%%%%%%%%%%%%%%%

\begin{applicationbox}

\textbf{Engineering Design.}
Designers balance weight, cost, reliability, performance, energy use, and
safety.

\medskip

\textbf{Finance.}
Portfolio decisions balance expected return, variance, liquidity, and
tail risk.

\medskip

\textbf{Transportation.}
Routing and network decisions may balance cost, travel time, emissions,
and service quality.

\medskip

\textbf{Supply Chains.}
Planning can simultaneously address cost, resilience, delivery time, and
environmental impact.

\medskip

\textbf{Healthcare.}
Resource-allocation decisions may balance cost, patient outcomes, access,
and waiting time.

\medskip

\textbf{Energy Systems.}
Power-system planning may balance cost, reliability, emissions, and
renewable penetration.

\medskip

\textbf{Environmental Planning.}
Land, water, and conservation decisions commonly require tradeoffs between
economic and ecological objectives.

\medskip

\textbf{Public Policy.}
Government decisions often involve several stakeholders with competing
objectives rather than one universally accepted scalar measure.

\end{applicationbox}

%%%%%%%%%%%%%%%%%%%%%%%%%%%%%%%%%%%%%%%%%%%%%%%%%%%%%%%%%%%%
\subsection{Exercises}
\label{subsec:multiobjective-exercises}
%%%%%%%%%%%%%%%%%%%%%%%%%%%%%%%%%%%%%%%%%%%%%%%%%%%%%%%%%%%%

\begin{exercise}[1]
Define multiobjective optimization.
\end{exercise}

\begin{exercise}[1]
Define Pareto dominance.
\end{exercise}

\begin{exercise}[1]
Define a Pareto-optimal solution.
\end{exercise}

\begin{exercise}[1]
Define the Pareto set.
\end{exercise}

\begin{exercise}[1]
Define the Pareto frontier.
\end{exercise}

\begin{exercise}[1]
Define weak Pareto optimality.
\end{exercise}

\begin{exercise}[1]
Define the ideal objective vector.
\end{exercise}

\begin{exercise}[1]
Define the nadir objective vector conceptually.
\end{exercise}

\begin{exercise}[1]
Define scalarization.
\end{exercise}

\begin{exercise}[1]
Define goal programming.
\end{exercise}

\begin{exercise}[2]
Assume both objectives are minimized. Determine whether

\[
(2,5)
\]

dominates

\[
(3,6).
\]
\end{exercise}

\begin{exercise}[2]
Determine whether either of the objective vectors

\[
(2,7)
\]

and

\[
(5,4)
\]

dominates the other.
\end{exercise}

\begin{exercise}[2]
From the objective vectors

\[
(2,8),
\quad
(3,5),
\quad
(4,7),
\quad
(6,3),
\]

identify all nondominated vectors.
\end{exercise}

\begin{exercise}[2]
Explain why the ideal objective vector may be infeasible.
\end{exercise}

\begin{exercise}[2]
Normalize an objective value using given ideal and nadir estimates.
\end{exercise}

\begin{exercise}[2]
Write a weighted-sum scalarization for three objectives.
\end{exercise}

\begin{exercise}[2]
Write an epsilon-constraint formulation with one primary objective and two
secondary objectives.
\end{exercise}

\begin{exercise}[2]
Write a lexicographic optimization sequence for three objectives.
\end{exercise}

\begin{exercise}[2]
Write goal-programming deviation equations for two target values.
\end{exercise}

\begin{exercise}[3]
For three feasible alternatives with two objectives, identify the Pareto
frontier.
\end{exercise}

\begin{exercise}[3]
Evaluate a weighted-sum objective using several different weight vectors.
\end{exercise}

\begin{exercise}[3]
Show that a strictly positive weighted-sum optimum is Pareto optimal.
\end{exercise}

\begin{exercise}[3]
Give a geometric example of an unsupported Pareto point.
\end{exercise}

\begin{exercise}[3]
Explain why weighted sums may fail on a nonconvex objective set.
\end{exercise}

\begin{exercise}[3]
Generate several Pareto solutions using the epsilon-constraint method.
\end{exercise}

\begin{exercise}[3]
Compare the solutions obtained from weighted-sum and epsilon-constraint
scalarizations.
\end{exercise}

\begin{exercise}[3]
Formulate a biobjective linear program.
\end{exercise}

\begin{exercise}[3]
Convert the previous model into a weighted linear program.
\end{exercise}

\begin{exercise}[3]
Convert the same model into an epsilon-constraint linear program.
\end{exercise}

\begin{exercise}[3]
Construct a goal-programming model with profit, overtime, and emissions
targets.
\end{exercise}

\begin{exercise}[3]
Construct a lexicographic model in which safety has higher priority than
cost.
\end{exercise}

\begin{exercise}[3]
Write a Chebyshev scalarization for two normalized objectives.
\end{exercise}

\begin{exercise}[3]
Convert a max-deviation objective into epigraph form.
\end{exercise}

\begin{exercise}[3]
Find the ideal point for a two-objective problem by solving the two
single-objective problems.
\end{exercise}

\begin{exercise}[3]
Explain how an inaccurate scaling range can distort weighted objectives.
\end{exercise}

\begin{exercise}[3]
For a differentiable unconstrained two-objective problem, write the
weighted first-order Pareto stationarity condition.
\end{exercise}

\begin{exercise}[3]
Write the constrained multiobjective KKT stationarity equation.
\end{exercise}

\begin{exercise}[3]
Explain the difference between a priori and a posteriori methods.
\end{exercise}

\begin{exercise}[3]
Explain the advantage of interactive multiobjective optimization.
\end{exercise}

\begin{exercise}[3]
Construct a mean--variance portfolio problem as a two-objective model.
\end{exercise}

\begin{exercise}[3]
Construct an engineering model minimizing both mass and cost.
\end{exercise}

\begin{exercise}[3]
Construct a transportation model minimizing cost and emissions.
\end{exercise}

\begin{exercise}[3]
Formulate a max--min fairness problem.
\end{exercise}

\begin{exercise}[3]
Explain why fairness and efficiency may conflict.
\end{exercise}

\begin{exercise}[3]
Explain what a knee point represents.
\end{exercise}

\begin{exercise}[3]
Explain why different decisions can map to the same Pareto objective
vector.
\end{exercise}

\begin{exercise}[3]
Explain why a Pareto frontier alone may not determine the final decision.
\end{exercise}

\begin{exercise}[4]
Study existence of Pareto-optimal solutions under compactness and
continuity assumptions.
\end{exercise}

\begin{exercise}[4]
Prove that every Pareto-optimal point is weakly Pareto optimal.
\end{exercise}

\begin{exercise}[4]
Give an example showing the converse fails.
\end{exercise}

\begin{exercise}[4]
Study scalarization theorems for convex multiobjective optimization.
\end{exercise}

\begin{exercise}[4]
Characterize supported and unsupported efficient points geometrically.
\end{exercise}

\begin{exercise}[4]
Study proper efficiency in the sense of Geoffrion.
\end{exercise}

\begin{exercise}[4]
Derive multiobjective KKT conditions using scalarization.
\end{exercise}

\begin{exercise}[4]
Study second-order conditions for local Pareto optimality.
\end{exercise}

\begin{exercise}[4]
Study achievement scalarizing functions.
\end{exercise}

\begin{exercise}[4]
Investigate reference-point methods.
\end{exercise}

\begin{exercise}[4]
Study Benson-type algorithms for multiobjective linear programming.
\end{exercise}

\begin{exercise}[4]
Study epsilon-constraint enumeration for discrete Pareto sets.
\end{exercise}

\begin{exercise}[4]
Investigate multiobjective branch-and-bound.
\end{exercise}

\begin{exercise}[4]
Study multiobjective dynamic programming.
\end{exercise}

\begin{exercise}[4]
Investigate evolutionary multiobjective optimization.
\end{exercise}

\begin{exercise}[4]
Study nondominated sorting algorithms.
\end{exercise}

\begin{exercise}[4]
Study the hypervolume indicator.
\end{exercise}

\begin{exercise}[4]
Investigate diversity metrics for approximate Pareto fronts.
\end{exercise}

\begin{exercise}[4]
Study robust multiobjective optimization.
\end{exercise}

\begin{exercise}[4]
Study stochastic multiobjective optimization.
\end{exercise}

\begin{exercise}[4]
Investigate preference elicitation for multiobjective decision support.
\end{exercise}

%%%%%%%%%%%%%%%%%%%%%%%%%%%%%%%%%%%%%%%%%%%%%%%%%%%%%%%%%%%%
\subsection{Conceptual Summary}
\label{subsec:multiobjective-summary}
%%%%%%%%%%%%%%%%%%%%%%%%%%%%%%%%%%%%%%%%%%%%%%%%%%%%%%%%%%%%

A multiobjective optimization problem has vector objective

\[
\boxed{
\mathbf{f}(x)
=
\begin{pmatrix}
f_1(x)\\
\vdots\\
f_k(x)
\end{pmatrix}.
}
\]

A feasible point is Pareto optimal when no other feasible point improves
one objective without worsening another.

The Pareto set lies in decision space:

\[
\boxed{
\mathcal{P}
=
\{
x:
x\text{ is Pareto optimal}
\}.
}
\]

Its image in objective space is the Pareto frontier.

Important scalarization methods include:

\[
\boxed{
\text{weighted sums},
}
\]

\[
\boxed{
\text{epsilon constraints},
}
\]

\[
\boxed{
\text{goal programming},
}
\]

\[
\boxed{
\text{lexicographic optimization},
}
\]

and

\[
\boxed{
\text{reference-point methods}.
}
\]

Weighted sums solve

\[
\boxed{
\min_x
\sum_i
w_if_i(x).
}
\]

Epsilon constraints solve one objective while bounding the others.

The ideal vector contains each objective's individual optimum, but is not
generally feasible.

Multiobjective optimization therefore combines

\[
\boxed{
\text{optimization}
+
\text{tradeoff analysis}
+
\text{preference modeling}
+
\text{decision support}.
}
\]

%%%%%%%%%%%%%%%%%%%%%%%%%%%%%%%%%%%%%%%%%%%%%%%%%%%%%%%%%%%%
\subsection{Section Connections}
\label{subsec:multiobjective-connections}
%%%%%%%%%%%%%%%%%%%%%%%%%%%%%%%%%%%%%%%%%%%%%%%%%%%%%%%%%%%%

\chapterconnection{
Multiobjective Optimization extends the single-objective framework of
Sections~\ref{sec:optimization-fundamentals},
\ref{sec:convex-optimization}, and
\ref{sec:stochastic-optimization} to problems where several competing
performance criteria must be considered simultaneously. Pareto dominance
replaces ordinary scalar comparison, while weighted sums, epsilon
constraints, goal programming, lexicographic methods, and reference-point
methods provide mechanisms for selecting compromises. These ideas connect
directly with Decision Theory later in Chapter~10. The next section,
Network Optimization, develops optimization problems whose feasible
structures are paths, flows, cuts, and networks, including shortest paths,
maximum flow, minimum-cost flow, transportation, and network-design
models.
}

%%%%%%%%%%%%%%%%%%%%%%%%%%%%%%%%%%%%%%%%%%%%%%%%%%%%%%%%%%%%
% SECTION SOURCE TRACKING
%%%%%%%%%%%%%%%%%%%%%%%%%%%%%%%%%%%%%%%%%%%%%%%%%%%%%%%%%%%%

%------------------------------------------------------------
% 10.9 Multiobjective Optimization
%------------------------------------------------------------
%
% Source basis:
%   - Standard multiobjective optimization theory.
%   - Standard vector optimization and Pareto-efficiency concepts.
%   - Standard multiple-criteria decision-making methods.
%
% Used for:
%   - vector-valued objectives;
%   - conflicting objectives;
%   - decision and objective spaces;
%   - Pareto dominance;
%   - dominated and nondominated solutions;
%   - Pareto optimality;
%   - weak Pareto optimality;
%   - Pareto sets and Pareto frontiers;
%   - tradeoff interpretation;
%   - ideal, nadir, and utopia reference points;
%   - objective normalization;
%   - scalarization;
%   - weighted-sum methods;
%   - supported and unsupported Pareto points;
%   - epsilon-constraint methods;
%   - lexicographic optimization;
%   - goal programming;
%   - reference-point methods;
%   - compromise programming;
%   - Chebyshev scalarization;
%   - multiobjective LP and convex optimization;
%   - first-order Pareto conditions;
%   - multiobjective KKT conditions;
%   - proper efficiency;
%   - preference information;
%   - a priori, a posteriori, and interactive methods;
%   - Pareto-front generation;
%   - knee points;
%   - discrete multiobjective optimization;
%   - evolutionary multiobjective optimization preview;
%   - nondominated sorting;
%   - diversity and hypervolume;
%   - weight sensitivity;
%   - fairness;
%   - sustainability;
%   - portfolio and engineering applications;
%   - exercises.
%
% Notes:
%   - Network Optimization is developed next in Section 10.10.
%   - Decision Theory is developed in Section 10.14.
%   - Stochastic multiobjective models connect with Section 10.8.
%
%%%%%%%%%%%%%%%%%%%%%%%%%%%%%%%%%%%%%%%%%%%%%%%%%%%%%%%%%%%%